\section{Experiments}\label{sec:exp}

We integrate \mname{} as a plug-and-play process during inference with 7B SOTA Video-LLMs.
Following the setup in~\citet{yu2024frame}, we uniformly sample 128 candidate frames and apply \mname{} to select 8 frames as input to the Video-LLMs.
We conduct experiments using LMMs-Eval~\cite{lmms_eval2024} and VLMEvalKit~\cite{duan2024vlmevalkit} on three long video benchmarks.
Due to space limitations, additional experimental details and results are provided in the Appendix.
Code in \url{https://github.com/sunh-23/MDP3}.

\begin{figure*}[!t]
    \centering
    \begin{minipage}{0.64\textwidth}
     	
\setlength\tabcolsep{3.5pt}
\begin{tabular}{lcccccc}
    \toprule
		\multirow{2}{*}{\textbf{Model}} & 
		\multicolumn{5}{c}{\textbf{LongVideoBench$_{val}$}} \\

		\cmidrule(lr){2-6}
		\quad\ \textit{Video Duration} & 
			Overall  & 
			(8s, 15s] & 
			(15s, 1m] & 
			(3m, 10m] & 
			(15m, 60m] \\ 
  	\toprule
	VILA-V1.5       & 47.1\emptyfoot{+3.7} & 56.1\emptyfoot{+1.6} & 60.5\emptyfoot{+5.8} & 43.4\emptyfoot{+4.9} & 42.7\emptyfoot{+2.9} \\
	\rowcolor{gray!30}
	\rule{0pt}{2ex} \textbf{+\mname}
	                & 50.8\gainfoot{+3.7} & 57.7\gainfoot{+1.6} & 66.3\gainfoot{+5.8} & 48.3\gainfoot{+4.9} & 45.6\gainfoot{+2.9} \\
	\hdashline
	MiniCPM-V2.6    & 51.2\emptyfoot{+5.9} & 67.7\emptyfoot{-0.5} & 66.9\emptyfoot{+4.0} & 46.8\emptyfoot{+9.3} & 44.1\emptyfoot{+6.3} \\
	\rowcolor{gray!30}
	\rule{0pt}{2ex} \textbf{+\mname} 
	                & 57.1\gainfoot{+5.9} & 67.2\lossfoot{-0.5} & 70.9\gainfoot{+4.0} & 56.1\gainfoot{+9.3} & 50.4\gainfoot{+6.3} \\
	\hdashline
	LLaVA-OneVision & 54.2\emptyfoot{+4.8} & 68.3\emptyfoot{+2.1} & 66.9\emptyfoot{+6.4} & 52.4\emptyfoot{+5.4} & 46.8\emptyfoot{+5.0} \\
	\rowcolor{gray!30}
	\rule{0pt}{2ex} \textbf{+\mname}
	                & 59.0\gainfoot{+4.8} & 70.4\gainfoot{+2.1} & 73.3\gainfoot{+6.4} & 57.8\gainfoot{+5.4} & 51.8\gainfoot{+5.0} \\
	\bottomrule
	\end{tabular}

     	\captionof{table}{
		    Performance of \mname{} and its absence on LongVideoBench$_{val}$ for varying video durations across different baseline Video-LLMs.
		}
		\label{tab:lvbv}
    \end{minipage}\hfill
    \begin{minipage}{0.34\textwidth}
    	\centering
	   	\includegraphics[width=1.0\linewidth]{figures/lvbv_gain.pdf}
		\captionof{figure}{
    		Performance of baselines and the improvement of \mname{} on LongVideoBench$_{val}$ for different video durations.
    		}
	\label{fig:lvbv}
    \end{minipage}
\end{figure*}


\subsection{Experiments Protocol}
\textbf{Backbone Models \& Primary Baselines:}
Based on the Video-MME leaderboard~\cite{fu2024video}, 
we choose the 7B SOTA Video-LLMs LLaVA-OneVision-7B~\cite{li2024llava}, MiniCPM-V2.6-7B~\cite{yao2024minicpm} and Ovis2-8B~\cite{ovis2} as our primary baselines.
To compare with the recent frame selection approach~\cite{yu2024frame}, tailor-made for VILA~\cite{lin2024vila}, we also include VILA-V1.5-7B as the primary baseline.
We integrate \mname{} with these baselines and evaluate it against other recent Video-LLMs.

\textbf{Benchmarks:}
We evaluate \mname{} on three long video benchmarks:
1)~\textbf{Video-MME}~\cite{fu2024video}, comprising 2700 human-annotated question-answer pairs, with an average video duration of 17 minutes. 
2)~\textbf{MLVU}~\cite{MLVU}, a comprehensive benchmark for multi-task long video understanding, consisting of 2593 tasks from 9 categories, with an average video duration of 12 minutes.
3)~\textbf{LongVideoBench}~\cite{wu2024longvideobench}, for which we use the validation set without subtitles, denoted as LVB$_{val}$, which contains 1337 question-answer pairs and has an average video length of 12 minutes.



\subsection{Results and Analysis}


\subsubsection{Comparison with SOTA Methods} 

\cref{tab:main_results} presents a comparison of \mname{} while applied to VILA-V1.5~\cite{lin2024vila},
MiniCPM-V2.6~\cite{yao2024minicpm}, LLaVA-OneVision~\cite{li2024llava} and Ovis2~\cite{ovis2}.
%We compare \mname{} to recent Video-LLMs, including 
We compare \mname{} to recent Video-LLMs:
Video-LLaVA~\cite{lin2023video}, Qwen-VL-Chat~\cite{bai2023qwen}, VideoChat2~\cite{li2023videochat}, Chat-UniVi-V1.5~\cite{jin2024chat}, VideoLLaMA2~\cite{cheng2024videollama2}, LLaVA-NEXT-QW2~\cite{liu2024llavanext}, LongVILA~\cite{xue2024longvila}, LongVA~\cite{zhang2024longva}, and Video-XL~\cite{shu2024video}. 
%
The number of input frames for the LLMs is shown in \cref{tab:main_results}. 
Specifically, Video-XL processes 128 frames for Video-MME and 256 frames for MLVU, while LLaVA-OneVision$^*$ reports the official results using a well-tuned frame count.	
%
% 
Compared to the baselines, \mname{} substantially improves the SOTA performance across all three evaluated benchmarks.
The main results presented in \cref{tab:main_results} highlight several key observations:
\begin{enumerate}
	\item Our frame selection method, \mname{}, empowers Ovis2 to achieve the highest overall performance across various benchmarks, confirming the effectiveness of \mname{}.
	\item When \mname{} is applied to the various primary Video-LLMs,
		it consistently outperforms uniform frame sampling.
		This validates that \mname{} is effective in frame selection and generalizes well across various Video-LLMs.
	\item Notably, LLaVA-OneVision with \mname{} selecting 8 frames outperforms LLaVA-OneVision$^*$, which uses a well-tuned frame count.
		This demonstrates that selecting key frames is more effective and efficient than simply increasing the number of frames, a finding supported by comparisons with LongVA, LongVILA, and Video-XL.
	\item Compared to \textsc{Frame-Voyager}, a custom-built method for VILA-V1.5 trained on extensive data, \mname{} shows superior performance. 
	This result validates that even a training-free, model-agnostic method can effectively achieve frame selection by reusing pretrained VLMs.
\end{enumerate}


\subsubsection{Results across Various Selection Sizes}\label{sssec:ravss}

\cref{fig:gain_on_videomme} presents the results of LLaVA-OneVision on Video-MME (without subtitles), comparing performance with and without applying \mname{} for key frame selection. 
There are four noteworthy insights:
\begin{enumerate}
	\item Regardless of whether \mname{} is applied, the law of diminishing marginal utility is evident as the number of frames fed into Video-LLMs increases.
		This shows that as the frame count rises, the visual understanding capacity of Video-LLMs quickly reaches its inherent upper bound, emphasizing the necessity for effective frame selection within limited capacity.
	\item When \mname{} is applied, the performance of LLaVA-OneVision consistently improves across various selection sizes, demonstrating the stability and generalizability of \mname{} under different selection size settings.
	\item The improvement of \mname{} is more pronounced when selecting fewer frames rather than more. 
		Notably, selecting 8 frames outperforms the best baseline of uniformly sampling 32 frames.
		Selecting just 2 frames (1/16 of the best baseline) achieves 90.8\% of the top performance, while 4 frames (1/8) reach 96.3\%.
		This shows that \mname{} effectively reduces the number of frames required for input, thereby enhancing inference efficiency.
	\item Notably, performance with 64 frames is lower than with 32 frames, whether or not \mname{} is used.
	This may be because LLaVA-OneVision was only trained on 32 frames, leading to potential performance drops when more frames are input.
	This indicates that current Video-LLMs often struggle to generalize beyond their training frame counts.
	Thus, effective frame selection methods like \mname{} are crucial for improving performance by focusing on frame selection during inference and avoiding the high costs of increasing frame counts during training.
\end{enumerate}


\subsubsection{Results across Various Video Durations}


We present the performance of baseline Video-LLMs and the improvements by \mname{} on LongVideoBench$_{val}$ for various video durations, as shown in \cref{tab:lvbv} and \cref{fig:lvbv}.
The results show that \mname{} substantially improves baseline performance for most video durations. 
However, for videos lasting 8 to 15 seconds, the improvements are smaller, and MiniCPM-V2.6 even shows a slight decline.
This is because uniformly sampling 8 frames in this range results in about 1 fps, making frame selection less impactful for very short videos.
This highlights that frame selection is more crucial for long videos than short ones. 
Indeed, for effectiveness and efficiency in video understanding, focusing on long videos provides a higher return on investment~(ROI).



\subsubsection{Ablation Study}

\begin{table}[t]
	\centering
	\setlength\tabcolsep{2.0pt}
	\begin{tabular}{lccccccc}
		\toprule
	                  & \textit{Qry.} & \textit{Div.} & \textit{Seq.} & \textit{HDS} & \textit{TFMA} & Acc. \\ 
		\midrule
	VILA-8B (+ Uniform)    & \nocheckmark  & \nocheckmark & \nocheckmark  & \nocheckmark & \checkmark & 47.5 \\
	\ + RGB Histogram      & \nocheckmark  & \nocheckmark & \nocheckmark  & \nocheckmark & \checkmark & 45.9 \\
	\ + Edges Change       & \nocheckmark  & \nocheckmark & \nocheckmark  & \nocheckmark & \checkmark & 47.3 \\
	\ + Optical Flow       & \nocheckmark  & \nocheckmark & \nocheckmark  & \nocheckmark & \checkmark & 46.7 \\
	\ + Katna              & \nocheckmark  & \nocheckmark & \nocheckmark  & \nocheckmark & \checkmark & 45.7 \\
	\ + \textsc{SeViLA}    
						  & \semicheckmark & \semicheckmark & \semicheckmark  & \nocheckmark & \nocheckmark &  49.0 \\
	\ + \textsc{Frame-Voyager}           
						   & \checkmark    & \semicheckmark  & \semicheckmark & \nocheckmark & \nocheckmark & 50.5 \\
		\hdashline
	\ + SigLIP 	           & \checkmark    & \nocheckmark & \nocheckmark  & \nocheckmark & \checkmark & 50.6 \\
	\ + \mname{} w. MGK    & \nocheckmark  & \checkmark   & \checkmark  & \checkmark & \checkmark & 48.9 \\
	\ + DPP w. CMGK  	   & \checkmark    & \checkmark   & \nocheckmark  & \checkmark & \checkmark & 51.8 \\
	\ + \mname{} w. cos. sim. & \checkmark & \checkmark   & \checkmark  & \nocheckmark & \checkmark & 50.2 \\
		\rowcolor{gray!30}
	\ \textbf{+ \mname}    & \checkmark    & \checkmark   & \checkmark    & \checkmark & \checkmark & \textbf{53.3} \\	
		\bottomrule
	\end{tabular}
	\caption{
		Results (\%) of various frame selection methods and ablation variants of \mname{} on Video-MME ($wo\ subs$), based on VILA. 
		Indicators: query relevance~(\textit{Qry}), list-wise diversity~(\textit{Div}), list-wise sequentiality~(\textit{Seq}), high-dimensional RKHS similarity (\textit{HDS}), and a tuning-free, model-agnostic design~(\textit{TFMA}).
%		\textit{Qry}: whether considering query relevance;
%		\textit{Div}: whether considering list-wise diversity;
%		\textit{Seq}: whether considering list-wise sequentiality;
%		\textit{HDS}: whether employ high-dimensional similarity in RKHS;
%		\textit{TFMA}: whether tuning-free and model-agnostic for Video-LLMs.
	}
	\label{tab:abl}
\end{table}



\begin{figure}[t]
	\centering
	\includegraphics[width=0.8\linewidth]{figures/frame_sel_comparison.pdf}
	\caption{
	    Accuracies (\%) of various ablation variants of \mname{} applied with MiniCPM-V2.6 on Video-MME (without subtitles).
	}
   \label{fig:abl}
\end{figure}


We conducted an ablation study of \mname{}, with the results shown in \cref{tab:abl} and \cref{fig:abl}.
The study considered variations such as 
directly using SigLIP to select the top-$k$ similar frames based on the query, 
using standard DPP with CMGK while disregarding sequentiality,
and applying \mname{} with similarity matrix modifications such as MGK and cosine similarity.
%
We also compared \mname{} with traditional rule-based shot boundary detection methods,
such as RGB Histogram~\cite{sheena2015key}, Edges Change Ratio~\cite{nasreen2013key}, and Optical Flow~\cite{wolf1996key}, which select frames based on significant transitions between adjacent frames. 
%
These methods do not account for the current query or the list-wise relationship among selected frames.
The cluster-based method Katna~\cite{keplerlab_katna} only considers pair-wise diversity, overlooking query relevance, list-wise diversity, and sequentiality. 
%
\textsc{Frame-Voyager}~~\cite{yu2024frame} and \textsc{SeViLA}~~\cite{yu2023self} are recent end-to-end frame-selection methods designed for a specific model.
The \semicheckmark{} symbol indicates that some factors, while not explicitly addressed, may be implicitly captured.
Training on VILA's hidden states with a list-based loss enables \textsc{Frame-Voyager} to potentially capture list-wise information, while position embeddings may address sequentiality.


\cref{tab:abl} shows that the key principles of query relevance, list-wise diversity, and sequentiality are crucial for frame selection.
The complete \mname{} is the first to incorporate all these principles and achieves the best performance in the ablation study.
%
A comparison of the ablation variants of \mname{} further confirms that 
all components are effective and indispensable, including high-dimensional similarity in RKHS, list-wise diversity with DPP, and sequentiality with MDP.
%
\cref{fig:abl} further illustrates the ablation results across various video durations. 
While \mname{} outperforms other methods overall, it performs slightly worse than DPP for short videos.
Standard DPP accounts for list-wise diversity but lacks sequentiality.
However, the performance of standard DPP deteriorates significantly with longer videos.
%
These observations reveal two conclusions:
\begin{enumerate}
    \item Sequentiality is more critical for long videos than for short ones, and it should be prioritized when selecting frames from longer videos.
    \item Standard DPP cannot handle sequentiality, which becomes a drawback as video length increases.  Therefore, sequentiality is essential for frame selection and \mname{}.
\end{enumerate}
















